\documentclass[paper=a4, fontsize=11pt]{scrartcl} % A4 paper and 11pt font size

\newcommand{\assignment}{5}
\newcommand{\duedate}{March 15, 2023}
\usepackage[top=1in, bottom=1.5in, left=1in, right=1in]{geometry}
\usepackage{fancyhdr} % Required for custom headers
\usepackage{lastpage} % Required to determine the last page for the footer
\usepackage{extramarks} % Required for headers and footers
\usepackage[usenames,dvipsnames]{color} % Required for custom colors
\usepackage{graphicx} % Required to insert images
\usepackage{listings} % Required for insertion of code
\usepackage{courier} % Required for the courier font
\usepackage{amsmath}
\usepackage[super]{nth}
\usepackage{booktabs}
\usepackage[usenames,dvipsnames]{xcolor}
\usepackage{tcolorbox}
\usepackage{tabularx}
\usepackage{array}
\usepackage{colortbl}

%\usepackage[T1]{fontenc} % Use 8-bit encoding that has 256 glyphs
%\usepackage{fourier} % Use the Adobe Utopia font for the document - comment this line to return to the LaTeX default
\usepackage[english]{babel} % English language/hyphenation
\usepackage{amsmath,amsfonts,amsthm} % Math packages
\usepackage{graphicx}

\usepackage{hyperref}
\hypersetup{
  colorlinks   = true, %Colours links instead of ugly boxes
  urlcolor     = blue, %Colour for external hyperlinks
  linkcolor    = blue, %Colour of internal links
  citecolor   = red %Colour of citations
}

\usepackage{fancyhdr} % Custom headers and footers
\pagestyle{fancyplain} % Makes all pages in the document conform to the custom headers and footers
\fancyhead{} % No page header - if you want one, create it in the same way as the footers below
\fancyfoot[L]{} % Empty left footer
\fancyfoot[C]{} % Empty center footer
\fancyfoot[R]{\thepage} % Page numbering for right footer
\renewcommand{\headrulewidth}{0pt} % Remove header underlines
\renewcommand{\footrulewidth}{0pt} % Remove footer underlines
\setlength{\headheight}{13.6pt} % Customize the height of the header
\newcommand{\ts}{\textsuperscript}

\numberwithin{equation}{section} % Number equations within sections (i.e. 1.1, 1.2, 2.1, 2.2 instead of 1, 2, 3, 4)
\numberwithin{figure}{section} % Number figures within sections (i.e. 1.1, 1.2, 2.1, 2.2 instead of 1, 2, 3, 4)
\numberwithin{table}{section} % Number tables within sections (i.e. 1.1, 1.2, 2.1, 2.2 instead of 1, 2, 3, 4)

\setlength\parindent{0pt} % Removes all indentation from paragraphs - comment this line for an assignment with lots of text

% Default fixed font does not support bold face
\DeclareFixedFont{\ttb}{T1}{txtt}{bx}{n}{8} % for bold
\DeclareFixedFont{\ttm}{T1}{txtt}{m}{n}{8}  % for normal

%----------------------------------------------------------------------------------------
%	CODE BLOCKS
%----------------------------------------------------------------------------------------

\usepackage{adjustbox}
\usepackage{listings}
\usepackage{color}

\definecolor{dkgreen}{rgb}{0,0.6,0}
\definecolor{gray}{rgb}{0.5,0.5,0.5}
\definecolor{mauve}{rgb}{0.58,0,0.82}

\lstdefinelanguage{Dockerfile}
{
  morekeywords={FROM, RUN, CMD, LABEL, MAINTAINER, EXPOSE, ENV, ADD, COPY,
    ENTRYPOINT, VOLUME, USER, WORKDIR, ARG, ONBUILD, STOPSIGNAL, HEALTHCHECK,
    SHELL},
  morecomment=[l]{\#},
  morestring=[b]"
}

\lstset{
    columns=flexible,
    aboveskip=5mm,
    belowskip=5mm,
    keepspaces=true,
    showstringspaces=false,
    basicstyle=\ttfamily,
    commentstyle=\color{gray},
    keywordstyle=\color{purple},
    stringstyle=\color{green}
}



%----------------------------------------------------------------------------------------
%	TITLE SECTION
%----------------------------------------------------------------------------------------

\usepackage{eso-pic}
% \usepackage[demo]{graphicx}
\newcommand\AtPageUpperRight[1]{\AtPageUpperLeft{%
   \makebox[\paperwidth][r]{#1}}}

\newcommand{\horrule}[1]{\rule{\linewidth}{#1}} % Create horizontal rule command with 1 argument of height

\title{	
\normalfont \normalsize
\textsc{Northeastern University,  Khoury College of Computer Science} \\ [25pt] % Your university, school and/or department name(s)
\horrule{0.5pt} \\[0.4cm] % Thin top horizontal rule
\huge CS 6220  Data Mining \textemdash~Assignment \assignment \\ % The assignment title
\Large \textbf{Due: \duedate (100 points)} % The assignment title
\horrule{2pt} \\[0.5cm] % Thick bottom horizontal rule
}

% Original in the document
\AddToShipoutPictureBG*{%
  \AtPageUpperRight{\raisebox{-\height}{\includegraphics[width=3cm]{include/logo}}}}

\author{
    \textbf{YOUR NAME} \\ 
    \textbf{YOUR GIT USERNAME} \\ 
    \textbf{YOUR E-MAIL}
}% INFORMATION

\date{} % Today's date or a custom date

\author{
    \textbf{YOUR NAME} \\ 
    \textbf{YOUR GIT USERNAME} \\ 
    \textbf{YOUR E-MAIL}
}% INFORMATION

\begin{document}

\maketitle % Print the title

{\huge \textbf{Na\"ive Bayes, Bayes Rules}} \\

The original performance of \href{https://www.ncbi.nlm.nih.gov/pmc/articles/PMC6339026/}{acoustic classification for Parkinsons Disease} leverages speech recordings from controlled subject responses from variety of questions. The task in the competition was to detect whether or not a person $X$ had Parkinsons disease from a sampling of data. As of 2018, the state of the art classifiers have achieved 90\% correct classification on a held out dataset, both for subjects who had Parkinsons and those who did not (at equal rates). So, when classifier $Y$ sees person $X$, it works correctly 90\% of the time. \\

{\Large \textbf{Question 1} [30 pts total]} \\
\\
Let's say that we run a clinic. This clinic leverages this classifier, which has 90\% accuracy. Also, let us say that we know that our current patient load is that 10\% of the population have Parkinsons and 90\% of the population do not. Let's also say that we're seeing patient $X$, and the classification algorithm has detected that they have Parkinson's disease. What's the probability that indeed $X$ has Parkinson's disease? \\

Come up with the numerical solution, and show your written work.

\vspace{1.5cm}
{\huge \textbf{Gradient Descent - Logistic Regression}}
\vspace{3mm}

The sinking of the RMS Titanic is the most notorious shipwrecks in history, sinking in the early hours of April 15, 1912. During her maiden voyage, the Titanic was widely considered “unsinkable”, but ultimately sank after colliding with an iceberg. In accordance with existing practice, Titanic's lifeboats did not account for all passengers onboard, and of the 2,224 passnegers, 1502 of them perished, most of whom were still on board when the ship sank. While seemingly random, it seems that some groups of people were more likely to survive than others. 

In this homework question, which is mirrored on \href{https://www.kaggle.com/competitions/titanic/overview}{Kaggle} (though the data samples are different), we ask you to build a predictive model that answers the question: “what sorts of people were more likely to survive?” using passenger data (i.e., name, age, gender, socio-economic class, etc). You can try getting hints on what features predicted survival from the \href{https://en.wikipedia.org/wiki/Sinking_of_the_Titanic#Casualties_and_survivors}{Wikipedia page}. \\

For \textbf{Question 2b.)}, we'll be using logistic regression with binary cross-entropy cost function appearing in the following form: \\

\begin{equation}
    \mathcal{L}(W, b) = -\sum_i y_i \log{ h_{W,b}(\textbf{x}) } + (1 - y_i) \log{( 1 - h_{W,b}(\textbf{x})) } \nonumber
\end{equation}
where 
\begin{equation}
h_{W,b}(\textbf{x}) = \sigma \left( W^T\textbf{x} + b \right) \nonumber
\end{equation}
and
\begin{equation}
\sigma(z) = \frac{1}{1 + \exp(-z)} \nonumber
\end{equation} \\
\\

\vspace{1cm}
{\Large \textbf{Question 2a.)} [20 pts]} \\

Before solving anything with Python, prove (on paper) that the derivative of the sigmoid function is 

\begin{equation}
    \nabla_z \sigma(z) = \sigma(z) \left(1 - \sigma(z)\right) \nonumber
\end{equation}

\vspace{1cm}
{\Large \textbf{Question 2b.)} [30 pts]} \\

We'll explore logistic regression with data that originated from the Kaggle site (but it is \emph{not} that exact data). Go ahead and download the data from \href{https://course.ccs.neu.edu/cs6220/homework-5/data/titanic/}{the homework 5 data folder}. Train a logistic regression predicting who would survive with \verb"titanic.train.csv". Test to see your accuracy on \verb"titanic.test.csv", where accuracy is defined by:

\begin{equation}
    \text{acc} = \frac{\text{num correct}}{\text{total}}
\end{equation}

Make sure you're not including the \verb"survived" column in the data as one of your features (that'd be cheating!) If you've done the lab in class, then this simply becomes an exercise of data transformations and preprocessing. (Since there is only one label that you are predicting, you may also have to ensure that the dimensions are correct.) \\

Some tips and tricks:

\begin{itemize}
    \item For features with multiple values, try making one-hot encoded data. For example, number of siblings, you can try to have a five column array. \\
    
    \begin{tabular}{|c|c|c|c|c|}
        \hline
         No Siblings & 1 Sibling & 2 Siblings & 3 Siblings & 4+ Siblings  \\
         \hline
                0    &    1      &      0     &       0  &    0\\
        \hline
    \end{tabular} \\
    \item Play around with your learning rate. If your learning is too high or too low, then you won't converge to the right answer.
    \item You don't use all the features. Just because they're there, doesn't mean they're always useful. (I ignored using the names.)
    \item Make sure you normalize your features to have similar ranges. You needn't use the standard scalar (i.e., z-scores), but if you don't transform some features, the logistic regression will not work well.
    \item While Kaggle provides good descriptions of what the data is, do \emph{NOT} use the Kaggle dataset. Only use data provided via \href{https://course.ccs.neu.edu/cs6220/homework-5/}{course website}.
\end{itemize}

See if you can do better than I can with a \textbf{logistic regression} by adopting the code from the lab in lecture. On the test set, with that code, I reached 79.2\% accuracy at a threshold of 0.5, using just a few features. \\

For this question, submit a plot of your loss curve and print your overall accuracy on training and test sets. Include what features you used, what you did to transform them, and also include what you used as a learning rate.

\vspace{1cm}
{\Large \textbf{Question 2c.) }[20 pts] } \\

Using a toolbox (like Keras), try benchmarking your logistic regression and then adding a few more layers (i.e., a neural network). Diagram your neural network architecture and print out your overall \textbf{test set} accuracy. Make sure you do \emph{not} train on any samples from your test set.

\end{document}
